%% GERADO por gera_tab_inversoes.py -- nao editar a mao
\begin{table}[!ht]
\centering
\caption{Every operation--type cell below the 50 per cent chance floor, with its \emph{decisive} accuracy $W/(W+L)$, which leaves ties out of the denominator and therefore differs from the strict accuracy of the main tables wherever an operation ties. Significance is an exact sign-flip randomization with Holm across the 120 operation--type tests, reported for both nested units. Over the 20 broadcast segments, 22 of the 35 cells survive and $^{\dagger}$ marks the 13 that do not; over the 12 topic blocks, the more conservative unit, 11 survive and $^{\ddagger}$ marks the further 11 that fall. The 5 cells in which an operation is wrong on all one hundred pairs carry neither mark: they survive in both units. Counting ties inside the denominator instead of outside leaves the same 35 cells and the same 15 operations. Removing the two types aligned with an operation's own attribute, darkening and desaturation, leaves all 15 operations still inverting; removing white noise instead leaves 9. Dropping one segment at a time, all 15 keep inverting in every one of the 20 rounds. \textbf{Ties are unevenly spread}: 4 operations tie on at least one pair and the largest count is 81. Where ties are many, accuracy below 50 per cent does not by itself mean the degraded copy is preferred; the reading that does is correct against incorrect, which is what the sign-flip test compares. Dk = darkening, Mo = motion blur, Bl = Gaussian blur, H2 = H.264, Rs = rescaling, St = desaturation, Ns = white noise, Jp = JPEG.}
\label{tab:inversoes}
\footnotesize
\begin{tabular}{@{}>{\raggedright\arraybackslash}p{3.4cm}cc>{\raggedright\arraybackslash}p{5.6cm}@{}}
\toprule
\rowcolor{headerblue}
\hd{Operation} & \hd{Types} & \hd{Worst} & \hd{Cells below chance (\%)} \\
\midrule
Saliency area & 6 & \textcolor{accentred}{\textbf{10.0}} & Mo~29.0$^{\dagger}$, Bl~20.0$^{\dagger}$, Rs~21.0$^{\dagger}$, St~18.0$^{\ddagger}$, Ns~13.1$^{\ddagger}$, Jp~10.0$^{\ddagger}$ \\
\rowcolor{rowlight}Hue count (Ke) & 5 & \textcolor{accentred}{\textbf{4.2}} & Dk~21.7$^{\dagger}$, Mo~15.7$^{\dagger}$, Bl~12.3$^{\ddagger}$, Rs~14.1$^{\dagger}$, St~4.2$^{\ddagger}$ \\
Gradient orientation & 4 & \textcolor{accentred}{\textbf{0.0}} & Bl~26.0$^{\dagger}$, Rs~9.0, St~21.0$^{\dagger}$, Ns~0.0 \\
\rowcolor{rowlight}Gradient magnitude & 3 & \textcolor{accentred}{\textbf{0.0}} & St~23.0$^{\dagger}$, Ns~0.0, Jp~5.0 \\
Saturation & 3 & \textcolor{accentred}{\textbf{1.0}} & Dk~13.0$^{\ddagger}$, H2~15.0$^{\ddagger}$, Ns~1.0 \\
\rowcolor{rowlight}Entropy & 2 & \textcolor{accentred}{\textbf{0.0}} & H2~10.0$^{\ddagger}$, Ns~0.0 \\
Contrast (Ke) & 2 & \textcolor{accentred}{\textbf{10.3}} & St~31.6$^{\dagger}$, Ns~10.3$^{\ddagger}$ \\
\rowcolor{rowlight}Edge distribution & 2 & \textcolor{accentred}{\textbf{11.0}} & H2~11.0$^{\ddagger}$, Jp~31.0$^{\dagger}$ \\
Grey-level dispersion & 2 & \textcolor{accentred}{\textbf{22.0}} & St~22.0$^{\dagger}$, Ns~25.0$^{\dagger}$ \\
\rowcolor{rowlight}GLCM contrast & 1 & \textcolor{accentred}{\textbf{0.0}} & Ns~0.0 \\
Laplacian variance & 1 & \textcolor{accentred}{\textbf{0.0}} & Ns~0.0 \\
\rowcolor{rowlight}Gygli mean of three & 1 & \textcolor{accentred}{\textbf{3.0}} & Ns~3.0 \\
Colourfulness (Hasler) & 1 & \textcolor{accentred}{\textbf{3.0}} & Ns~3.0 \\
\rowcolor{rowlight}Colourfulness (Ke) & 1 & \textcolor{accentred}{\textbf{3.0}} & Ns~3.0 \\
Srinivas composite & 1 & 29.0 & Ns~29.0$^{\ddagger}$ \\
\bottomrule
\end{tabular}
\end{table}
